\documentclass{situationist-preprint}

\title{The Conditions of the Public Sphere}
\msid{POL-20-a}
\msfamily{POL-20}
\msdiscipline{Political Philosophy}
\mstype{Principal manuscript}
\mskeywords{public sphere; deliberative democracy; Habermas; equal access; media; fair value of political liberty; attention}

\begin{document}
\maketitle

\begin{abstract}
Is a public sphere possible? Habermas described the public sphere as the domain in which private persons come together to form opinion on matters of common concern through the exchange of reasons, and traced its decline under mass media and organized interests. Skeptics have argued that it was never more than an ideal resting on exclusions, and that modern communication makes it impossible. We argue that a public sphere is possible, but only under institutional conditions that guarantee equal access to public debate. The argument has two parts. First, we show that equal access is necessary: a domain of opinion-formation in which access depends on unequally distributed resources does not satisfy the openness condition that distinguishes a public sphere from an arena of influence, and the markets in which public attention is allocated concentrate access without limit. Second, we show that equal access is achievable: the relevant inequality is one of opportunity to be heard, not of actual influence, and it can be secured by institutions such as public media, limits on political spending, rules for platforms of mass communication and public funding of political speech. We answer objections from Lippmann, Fraser and the marketplace of ideas.
\end{abstract}

\section{The problem}

\citet{Habermas1989} reconstructed the bourgeois public sphere of the eighteenth century, the coffee houses, salons and periodical press in which private persons debated public matters and claimed to hold state power to account by the standard of reasoned opinion. Its principles were that participants disregarded status, that any matter of common concern could be discussed, and that the public was in principle open to all. Habermas argued that the sphere was transformed in the nineteenth and twentieth centuries: the press became a commercial industry, organized interests and parties came to manufacture opinion rather than submit to it, and the public was reduced to an audience. In later work he described the public sphere as a network for communicating information and points of view, which filters and synthesizes them into public opinions that influence formal decision-making \citep{Habermas1996}.

Three skeptical lines of argument follow. The first, from \citet{Lippmann1922}, is that ordinary citizens cannot form competent opinions on public affairs, since they have neither the time nor the access to information; opinion is formed from stereotypes and can be managed. The second, from \citet{Fraser1990}, is that the bourgeois public sphere was constituted by the exclusion of women, workers and others, that the bracketing of status did not remove status inequalities but concealed them, and that a single comprehensive public sphere will always favor dominant groups. The third, prominent in recent discussion of digital media, is that communication fragmented into enclaves of the like-minded and organized by the competition for attention cannot sustain a common debate \citep{Sunstein2001}.

Our question is whether a public sphere is possible at all, under any feasible arrangement. Deliberative democrats require it, since on their view the legitimacy of democratic decisions depends on their emerging from free and equal public reasoning \citep{Cohen1989}. If it is impossible, deliberative democracy is a fiction.

\section{Framework}

\begin{definition}\label{d:sphere}
A \emph{public sphere} is a domain of communication in which (a) opinions on matters of common concern are formed (b) through the exchange of reasons, (c) which is open to all who are subject to the decisions at issue, and (d) whose results bear on the exercise of political power.
\end{definition}

The definition follows \citet{Habermas1996} in its four elements. Condition (c) is the one at issue. Openness cannot mean merely the absence of legal barriers to speech, since a domain in which everyone may speak but only some can be heard is not open in any sense that matters for the formation of opinion. We therefore make it precise.

\begin{definition}\label{d:access}
A citizen has \emph{access} to public debate to the extent that she has the opportunity to bring a view and the reasons for it to the attention of a substantial part of the public. \emph{Equal access} obtains when this opportunity does not depend on characteristics, such as wealth, social position or ownership of means of communication, that are unequally distributed for reasons unrelated to the merits of the view.
\end{definition}

Equal access is a condition on opportunity, not on outcome. It does not require that every view receive equal attention, which would be neither possible nor desirable, since views gain attention partly by being better supported. It requires that the differences in attention be explained by what is said and how it is received, rather than by who can pay.

\begin{principle}[Openness]\label{p:openness}
A domain of communication satisfies condition (c) of Definition~\ref{d:sphere} only to the extent that its participants have equal access to it.
\end{principle}

The idea is Rawls's guarantee of the fair value of the political liberties: formal rights of speech and participation are not enough if those with greater means can control the course of public debate \citep[\S36]{Rawls1971}. A domain in which access is proportioned to wealth is an arena in which the wealthy address an audience, not a public in which citizens reason together.

\begin{assumption}[Concentration]\label{a:concentration}
In the absence of institutions that counteract it, access to public attention tends to concentrate: those who already command attention or resources acquire more of it, since attention is scarce, visibility attracts further visibility, and the means of reaching a mass audience can be bought.
\end{assumption}

Assumption~\ref{a:concentration} is an empirical generalization supported by the history of the press that Habermas recounts, by the economics of broadcast and print media, and by the cumulative-advantage dynamics of networked communication, in which the distribution of audiences across sources is highly skewed. It does not say that concentration is inevitable, only that it is the tendency of an unregulated market for attention.

\section{The argument}

\begin{proposition}[Necessity]\label{p:necessity}
A public sphere exists only under institutional conditions that guarantee equal access to public debate.
\end{proposition}

\begin{proof}
By Principle~\ref{p:openness}, a domain of communication is a public sphere only to the extent that its participants have equal access. By Assumption~\ref{a:concentration}, access in the absence of counteracting institutions tends to concentrate, and so to depend on resources unequally distributed for reasons unrelated to the merits of views. Without such institutions, then, the domain tends away from equal access and so away from being a public sphere. Since the tendency is continuous, a domain that happens to be open at one time will not remain so without institutions that sustain its openness; and such institutions are what we mean by conditions that guarantee equal access.
\end{proof}

\begin{proposition}[Feasibility]\label{p:feasibility}
Institutional conditions that guarantee equal access to public debate can be realized.
\end{proposition}

\begin{proof}
It suffices to show that each of the main channels through which unequal resources become unequal access can be counteracted by an institution of a known type, without suppressing speech. (i) Ownership of the means of mass communication: public service media under independent governance, together with rules against concentration of ownership, provide channels whose access does not depend on ownership. (ii) Spending on political communication: limits on campaign and issue spending, disclosure requirements and public funding of parties and candidates reduce the conversion of wealth into reach. (iii) Control of digital platforms: transparency and non-discrimination rules for platforms that function as common channels of public communication prevent their owners or highest bidders from determining what reaches the public. (iv) The capacities required for participation: public education, translation and accessible public information distribute the competence to take part. Each institution has been realized in some democracy, and none requires that any view be excluded. Their joint realization is therefore feasible, and they jointly address the channels through which Assumption~\ref{a:concentration} operates.
\end{proof}

\begin{theorem}\label{t:main}
A public sphere is possible, but only under institutional conditions that guarantee equal access to public debate.
\end{theorem}

\begin{proof}
The \enquote{only under} clause is Proposition~\ref{p:necessity}. Possibility follows from Proposition~\ref{p:feasibility}: where the institutions are in place, the openness condition is satisfied to the extent required, and conditions (a), (b) and (d) of Definition~\ref{d:sphere} are not affected by the institutions except to the extent of being made easier to satisfy.
\end{proof}

The theorem places the public sphere among institutional achievements rather than among spontaneous social formations. The eighteenth-century public sphere that Habermas described was, on this view, not a lost natural condition but a partial and unstable approximation, made possible by a small and homogeneous reading public and destroyed by the mass market that extended it.

\section{Objections}

\begin{objection}[Lippmann]
The obstacle to a public sphere is not unequal access but the limited competence of citizens. Even with equal access, the public cannot form reasonable opinions on complex questions \citep{Lippmann1922}.
\end{objection}

\begin{reply}
\citet{Dewey1927} answered that the competence of the public is a function of the conditions of communication: a public that is poorly informed has been poorly served by the channels through which it learns. Experiments in structured deliberation show that ordinary citizens given balanced information and time to deliberate change their views in the direction of the evidence \citep{Fishkin2009}. Competence is therefore partly an effect of access, and condition (iv) of Proposition~\ref{p:feasibility} addresses it. Nor does a public sphere require that every citizen be competent on every question; it requires that the opinion formed through it be responsive to reasons, which a division of epistemic labor allows.
\end{reply}

\begin{objection}[Fraser]
A single public sphere favors dominant groups even under formal equality. Subordinated groups need counterpublics of their own, in which they can formulate their interests before entering a wider debate \citep{Fraser1990}.
\end{objection}

\begin{reply}
We agree, and the point supports the theorem. Fraser's argument is that formal openness conceals inequality of access; that is Principle~\ref{p:openness}. Subaltern counterpublics are among the institutions through which equal access is realized, since they give groups the opportunity to develop views and reasons that the dominant channels would not carry. The public sphere of Definition~\ref{d:sphere} need not be a single forum; it can be a network of publics connected so that views formed in one can reach the others \citep{Mansbridge2012}. What it cannot be is a network in which some publics are unable to reach the rest.
\end{reply}

\begin{objection}[The marketplace of ideas]
Regulating access to debate is a form of censorship. The best guarantee of open debate is the free competition of ideas, in which better views prevail \citep[ch.~II]{MillOL}.
\end{objection}

\begin{reply}
The argument for free competition among ideas assumes that ideas compete on their merits. Where reach can be bought, ideas compete partly on the resources of their sponsors, and the competition does not select for truth. None of the institutions of Proposition~\ref{p:feasibility} restricts what may be said; they restrict the conversion of money and ownership into reach. They are conditions of the marketplace of ideas in the sense Mill needed, not limitations of it.
\end{reply}

\begin{objection}[Fragmentation]
Even with equal access, digital communication fragments the public into enclaves that never encounter one another's reasons \citep{Sunstein2001}. A common public sphere is impossible however access is distributed.
\end{objection}

\begin{reply}
Fragmentation is partly a product of the unequal access the theorem addresses: platforms that sort users by engagement do so because attention is sold to the highest bidder, and the rules of condition (iii) can require that common channels expose users to a range of sources. But the deeper reply is that a public sphere does not require that everyone hear everything. It requires that views can reach the wider public when they warrant attention. A network of partially overlapping publics, connected by shared channels, satisfies Definition~\ref{d:sphere}; a set of sealed enclaves does not, and the difference is institutional.
\end{reply}

\section{Access and uptake}

Definition~\ref{d:access} concerns the opportunity to bring a view to the attention of the public. \citet{Young2000} distinguishes the \emph{external} exclusion of those kept out of a forum from the \emph{internal} exclusion of those who are present but whose contributions are dismissed, because their manner of speaking does not match the norms of the dominant group or because they are assigned less credibility. \citet{Fricker2007} analyzes the second phenomenon as testimonial injustice: a speaker receives less credibility than her evidence warrants because of prejudice about a group to which she belongs. If internal exclusion is pervasive, one might think that equal access in our sense is too weak to secure a public sphere, since a view can reach the public and still not be heard.

We accept that internal exclusion is a defect of a public sphere, and that equal access does not by itself remove it. But it does not show that the condition we state is insufficient in the sense the theorem requires. Credibility deficits are attached to speakers and groups, and they are reduced by exposure: views that reach the public repeatedly, from speakers of many kinds, change the expectations against which credibility is assigned. Access is therefore the condition under which internal exclusion can be corrected in public, while without access it cannot even be contested. Institutions can also address uptake directly, by norms of moderation in public forums, by the design of deliberative bodies so that every participant speaks, and by the presence in public media of speakers from groups that the dominant channels neglect. These are refinements of the institutions of Proposition~\ref{p:feasibility}, not a further condition of a different kind.

The point bears on Habermas's own later diagnosis. In his reconsideration of the public sphere under digital conditions, \citet{Habermas2023} argues that platforms have dissolved the editorial filters that once organized a common public, so that the problem is no longer only that access is unequal but that the inclusive, structured debate on which deliberative politics relies is replaced by fragmented semi-publics. On our account, this is a change in the channels through which Assumption~\ref{a:concentration} operates, not a refutation of the possibility of a public sphere. The filtering function of editors can be supplied in other forms, by public media, by platform rules and by deliberative institutions, and whether it is supplied is a matter of institutional choice.

\section{Relations to other results}

The theorem bears on several neighboring questions. POL-04 argues that democracy is justified both by the equal standing it confers and by its tendency to good decisions; both depend on a public sphere, the first because equal standing in voting is hollow without equal standing in the debate that precedes it, the second because the competence on which majority rule depends is formed in public debate. POL-14 argues for a chamber chosen by lot; an allotted chamber is one institution through which views unrepresented in the dominant channels reach the legislature. And POL-02 holds that a state is legitimate when it can be justified to all reasonable persons subject to it; a public sphere with equal access is the setting in which such justification can be offered and tested.

\section{What the result does not establish}

The theorem does not establish that any existing democracy has a public sphere; on our argument most have one only partially, to the extent that their institutions secure equal access. It does not establish which particular institutions are best, only that institutions of the types listed are jointly sufficient for the access required. And it does not establish that deliberation in a public sphere reaches correct conclusions, only that it is open in the way the conception requires.

\section{Conclusion}

A public sphere is a domain in which citizens form opinion on common concerns through reasons, open to all who are subject to the decisions at issue. Openness requires equal access to debate, and access left to markets in attention concentrates. A public sphere therefore exists only where institutions secure equal access. Those institutions are of known and realized types, and they restrict the conversion of resources into reach, not what may be said. The public sphere is possible, as an institutional achievement.

\begin{thebibliography}{99}\small
\bibitem[Cohen(1989)]{Cohen1989} Cohen, J. (1989). Deliberation and democratic legitimacy. In A. Hamlin and P. Pettit (eds.), \emph{The Good Polity: Normative Analysis of the State}. Blackwell.
\bibitem[Dewey(1927)]{Dewey1927} Dewey, J. (1927). \emph{The Public and Its Problems}. Henry Holt.
\bibitem[Fishkin(2009)]{Fishkin2009} Fishkin, J. S. (2009). \emph{When the People Speak: Deliberative Democracy and Public Consultation}. Oxford University Press.
\bibitem[Fraser(1990)]{Fraser1990} Fraser, N. (1990). Rethinking the public sphere: A contribution to the critique of actually existing democracy. \emph{Social Text}, 25/26, 56--80.
\bibitem[Fricker(2007)]{Fricker2007} Fricker, M. (2007). \emph{Epistemic Injustice: Power and the Ethics of Knowing}. Oxford University Press.
\bibitem[Habermas(1989)]{Habermas1989} Habermas, J. (1989). \emph{The Structural Transformation of the Public Sphere} (T. Burger and F. Lawrence, trans.). MIT Press. (Original work published 1962.)
\bibitem[Habermas(1996)]{Habermas1996} Habermas, J. (1996). \emph{Between Facts and Norms} (W. Rehg, trans.). MIT Press. (Original work published 1992.)
\bibitem[Habermas(2023)]{Habermas2023} Habermas, J. (2023). \emph{A New Structural Transformation of the Public Sphere and Deliberative Politics} (C. Cronin, trans.). Polity.
\bibitem[Lippmann(1922)]{Lippmann1922} Lippmann, W. (1922). \emph{Public Opinion}. Harcourt, Brace.
\bibitem[Mansbridge et~al.(2012)]{Mansbridge2012} Mansbridge, J., Bohman, J., Chambers, S., Christiano, T., Fung, A., Parkinson, J., Thompson, D. F., and Warren, M. E. (2012). A systemic approach to deliberative democracy. In J. Parkinson and J. Mansbridge (eds.), \emph{Deliberative Systems}. Cambridge University Press.
\bibitem[Mill(1859)]{MillOL} Mill, J. S. (1859). \emph{On Liberty}. John W. Parker and Son.
\bibitem[Rawls(1971)]{Rawls1971} Rawls, J. (1971). \emph{A Theory of Justice}. Harvard University Press.
\bibitem[Sunstein(2001)]{Sunstein2001} Sunstein, C. R. (2001). \emph{Republic.com}. Princeton University Press.
\bibitem[Young(2000)]{Young2000} Young, I. M. (2000). \emph{Inclusion and Democracy}. Oxford University Press.
\end{thebibliography}

\end{document}
